\documentclass[UTF8,11pt]{ctexart}
\usepackage[a4paper,margin=24mm]{geometry}
\usepackage{amsmath,amssymb,amsthm,mathtools,mathrsfs}
\usepackage[all]{xy}
\usepackage{xurl,hyperref,enumitem}
\usepackage{tikz}
\usetikzlibrary{arrows.meta,cd,decorations.markings,calc,patterns}
\hypersetup{hidelinks,breaklinks=true}
\setlength{\emergencystretch}{3em}
\allowdisplaybreaks
\newtheorem*{theorem}{Theorem}
\newtheorem*{thm}{Theorem}
\newtheorem*{lemma}{Lemma}
\newtheorem*{proposition}{Proposition}
\newtheorem*{prop}{Proposition}
\newtheorem*{corollary}{Corollary}
\newtheorem*{cor}{Corollary}
\newtheorem*{claim}{Claim}
\theoremstyle{definition}
\newtheorem*{definition}{Definition}
\newtheorem*{defn}{Definition}
\newtheorem*{remark}{Remark}
\newtheorem*{example}{Example}
\newcommand{\R}{\mathbb R}
\newcommand{\C}{\mathbb C}
\newcommand{\Z}{\mathbb Z}
\newcommand{\Q}{\mathbb Q}
\newcommand{\N}{\mathbb N}
\newcommand{\F}{\mathbb F}
\newcommand{\D}{\mathbb D}
\newcommand{\bB}{\mathbb B}
\newcommand{\bC}{\mathbb C}
\newcommand{\bR}{\mathbb R}
\newcommand{\bZ}{\mathbb Z}
\newcommand{\bN}{\mathbb N}
\newcommand{\bQ}{\mathbb Q}
\newcommand{\bD}{\mathbb D}
\newcommand{\bF}{\mathbb F}
\newcommand{\CP}{\mathbb{CP}}
\newcommand{\RP}{\mathbb{RP}}
\newcommand{\sO}{\mathscr O}
\newcommand{\sI}{\mathscr I}
\newcommand{\sF}{\mathscr F}
\newcommand{\sG}{\mathscr G}
\newcommand{\sH}{\mathscr H}
\newcommand{\sR}{\mathscr R}
\newcommand{\sM}{\mathscr M}
\newcommand{\sV}{\mathscr V}
\newcommand{\sU}{\mathscr U}
\DeclareMathOperator{\Hom}{Hom}
\DeclareMathOperator{\Ext}{Ext}
\DeclareMathOperator{\Tor}{Tor}
\DeclareMathOperator{\Spec}{Spec}
\DeclareMathOperator{\Proj}{Proj}
\DeclareMathOperator{\supp}{supp}
\DeclareMathOperator{\im}{im}
\DeclareMathOperator{\id}{id}
\DeclareMathOperator{\Tr}{Tr}
\DeclareMathOperator{\tr}{tr}
\DeclareMathOperator{\rank}{rank}
\DeclareMathOperator{\ord}{ord}
\DeclareMathOperator{\dist}{dist}
\DeclareMathOperator{\End}{End}
\DeclareMathOperator{\Aut}{Aut}
\DeclareMathOperator{\Gal}{Gal}
\DeclareMathOperator{\vol}{vol}
\DeclareMathOperator{\Cl}{Cl}
\DeclareMathOperator{\coker}{coker}
\DeclareMathOperator{\codim}{codim}
\DeclareMathOperator{\divergence}{div}
\DeclareMathOperator{\Ric}{Ric}
\DeclareMathOperator{\grad}{grad}
\DeclareMathOperator{\Hess}{Hess}
\newcommand{\abs}[1]{\left\lvert#1\right\rvert}
\newcommand{\sabs}[1]{\lvert#1\rvert}
\newcommand{\babs}[1]{\bigl\lvert#1\bigr\rvert}
\newcommand{\Babs}[1]{\Bigl\lvert#1\Bigr\rvert}
\newcommand{\bbabs}[1]{\biggl\lvert#1\biggr\rvert}
\newcommand{\BBabs}[1]{\Biggl\lvert#1\Biggr\rvert}
\newcommand{\norm}[1]{\left\lVert#1\right\rVert}
\newcommand{\snorm}[1]{\lVert#1\rVert}
\newcommand{\bnorm}[1]{\bigl\lVert#1\bigr\rVert}
\newcommand{\Bnorm}[1]{\Bigl\lVert#1\Bigr\rVert}
\newcommand{\bbnorm}[1]{\biggl\lVert#1\biggr\rVert}
\newcommand{\BBnorm}[1]{\Biggl\lVert#1\Biggr\rVert}
\newcommand{\widebar}[1]{\overline{#1}}
\newcommand{\nicefrac}[2]{\frac{#1}{#2}}
\newcommand{\linnprod}[2]{\langle#1,#2\rangle}
\newcommand{\myindex}[1]{#1}
\newcommand{\glsadd}[1]{}
\newcommand{\avoidbreak}{}
\newcommand{\sourceref}[1]{\textnormal{[原文：\nolinkurl{#1}]}}
\newcommand{\thmref}[1]{\sourceref{#1}}
\newcommand{\lemmaref}[1]{\sourceref{#1}}
\newcommand{\propref}[1]{\sourceref{#1}}
\newcommand{\corref}[1]{\sourceref{#1}}
\newcommand{\defnref}[1]{\sourceref{#1}}
\newcommand{\exampleref}[1]{\sourceref{#1}}
\newcommand{\exerciseref}[1]{\sourceref{#1}}
\newcommand{\figureref}[1]{\sourceref{#1}}
\newcommand{\sectionref}[1]{\sourceref{#1}}
\newcommand{\appendixref}[1]{\sourceref{#1}}
\newcommand{\stacksref}[2]{\href{https://stacks.math.columbia.edu/tag/#1}{#2 [#1]}}

\begin{document}
\section*{042\quad 辛流形上相容几乎复结构的构造与道路连通性}
\subsection*{来源与计数目标}
Kartik Venkatram 与 Denis Auroux 合作整理的 MIT 18.966 \emph{Geometry of Manifolds} 人类课程讲义，Spring 2007。课程页面明确记载该署名关系；不是将学生笔记冒充授课者独著。Lecture 7 第1页 Definition 1 与第2页 Proposition 1、Remark、Definition 2；Lecture 8 第1页 Proposition 1、Corollary 1。获取日期：2026-09-28。
\par\url{https://ocw.mit.edu/courses/18-966-geometry-of-manifolds-spring-2007/9188ed65ddf5e1662307d486b2c0418f_lect07.pdf}
\par\url{https://ocw.mit.edu/courses/18-966-geometry-of-manifolds-spring-2007/ae6565bbfacae2e53361e06acc340213_lect08.pdf}
\par\url{https://ocw.mit.edu/courses/18-966-geometry-of-manifolds-spring-2007/pages/lecture-notes/}
\par\textbf{本文件只计一个问题：}构造与辛形式相容的几乎复结构，并证明其空间道路连通。
\subsection*{作者命题与人类证明}

\begin{definition}[Lecture 7, Definition 1]
A complex structure on a vector space $V$ is an endomorphism $J:V\to V$ s.t. $J^2=-I$. Thinking of this $J$ as multiplication by $i$ turns $V$ into a complex vector space, $(x+iy)v=xv+yJv$. If $V$ is a symplectic vector space with symplectic form $\Omega$, a complex structure is compatible if $G(u,v)=\Omega(u,Jv)$ is a positive symmetric inner product. Note that being symmetric is equivalent to $\Omega(Ju,Jv)=\Omega(u,v)$, and being positive is precisely $\Omega(u,Ju)>0$ $\forall u\ne0$.
\end{definition}
\begin{proposition}[Lecture 7, Proposition 1]
If $(V,\Omega)$ is a symplectic vector space, $\exists$ a compatible $J$. Moreover, given any positive inner product $\langle\cdot,\cdot\rangle$ on $V$, we can build an $\Omega$-compatible complex structure on $V$ canonically (though it has no direct relation to the given inner product).
\end{proposition}
\begin{proof}
For the first part, taking $J=J_0$ in a standard basis gives the desired endomorphism. For the second part, by the nondegeneracy of $\Omega$, we have isomorphisms $u\mapsto\Omega(u,\cdot)$ and $u\mapsto\langle u,\cdot\rangle$ from $V$ to $V^*$. We thus obtain an endomorphism $A=\langle\rangle^{-1}\circ\Omega$ s.t. $\Omega(u,v)=\langle Au,v\rangle$. $A$ is invertible and skew-symmetric w.r.t. $\langle\rangle$, i.e. $A^*=-A$ (since $\Omega(v,u)=\langle Av,u\rangle=\langle v,A^*u\rangle=\langle A^*u,v\rangle=-\Omega(u,v)=-\langle Au,v\rangle$). Thus, $AA^*=-A^2$ is symmetric and positive definite, therefore diagonalizable with real, strictly positive eigenvalues. This implies the existence of a square root $\sqrt{AA^*}$ $(=\operatorname{diag}(\sqrt{\lambda_i}))$, so define $J=(\sqrt{AA^*})^{-1}A$. (Note that the decomposition $A=\sqrt{AA^*}J$ gives a ``polar decomposition'' of $A$.) $A$ commutes with $\sqrt{AA^*}$: letting $V_i$ be the eigenspace of $AA^*$ with eigenvalue $\lambda_i$, or similarly that of $\sqrt{AA^*}$ with eigenvalue $\sqrt{\lambda_i}$, we find that,
\[\forall v\in V_i,\quad(AA^*)Av=-A^3v=A(AA^*)v=\lambda_iAv\ \Longrightarrow\ Av\in V_i.\]
So $J$ also commutes with $A$ and with $\sqrt{AA^*}$, and thus is skew-symmetric
\[J^*=A^*(\sqrt{AA^*})^{-1}=-A(\sqrt{AA^*})^{-1}=-J\]
and orthogonal
\[J^*J=A^*(\sqrt{AA^*})^{-1}(\sqrt{AA^*})^{-1}A=\id.\]
In particular, $J^2=-J^*J=-\id$. For compatibility, note that
\begin{align*}
\Omega(Ju,Jv)&=\langle AJu,Jv\rangle=\langle JAu,Jv\rangle=\langle Au,v\rangle=\Omega(u,v),\\
\Omega(u,Ju)&=\langle Au,Ju\rangle=\langle-JAu,u\rangle
=\langle-(\sqrt{AA^*})^{-1}AAu,u\rangle\\
&=\langle(\sqrt{AA^*})^{-1}(AA^*)u,u\rangle=\langle(\sqrt{AA^*})u,u\rangle>0,
\end{align*}
thus completing the proof.
\end{proof}
\begin{remark}
Note that $G(u,v)=\Omega(u,Jv)=\langle\sqrt{AA^*}u,v\rangle$, so if $\langle\cdot,\cdot\rangle$ was already compatible with $\Omega$, then $AA^*=I$, $J=A$, $G=\langle\cdot,\cdot\rangle$.
\end{remark}
\begin{definition}[Lecture 7, Definition 2]
An almost-complex structure on a manifold $M$ is $J\in\End(TM)$ s.t. $J^2=-I$ (i.e. $\forall x\in M$, $J_x$ is a complex structure on $T_xM$). If $M=(M,\omega)$ is a symplectic manifold, $J$ is compatible if $\forall x\in M$, $J_x$ is $\omega_x$-compatible, with associated Riemannian metric $g_x(u_x,v_x)=\omega_x(u_x,J_xv_x)$. We say that $(\omega,g,J)$ is a compatible triple, with any two determining the third.
\end{definition}
\begin{proposition}[Lecture 8, Proposition 1]
For $(M,\omega)$ a symplectic manifold with Riemannian metric $g$, $\exists$ a canonical almost complex structure $J$ compatible with $\omega$.
\end{proposition}
\begin{proof}[Idea]
Do polar decomposition on every tangent space.
\end{proof}
\begin{corollary}[Lecture 8, Corollary 1]
Any symplectic manifold has compatible almost-complex structures, and the space of such structures is path connected.
\end{corollary}
\begin{proof}
For the first part, using a partition of unity gives a Riemannian metric, so the rest follows from the proposition. For the second part, given $J_0,J_1$, let $g_i=\omega(\cdot,J_i\cdot)$ for $i=0,1$ and set $g_t=(1-t)g_0+tg_1$. Each of these (for $t\in[0,1]$) is a metric, and gives an $\omega$-compatible $\widetilde J_t$ by polar decomposition, with $\widetilde J_0=J_0$ and $\widetilde J_1=J_1$.
\end{proof}

\subsection*{整理说明（非作者正文）}
允许的先修是有限维谱定理、正定平方根的光滑依赖、向量丛与分割统一。没有把相容复结构存在性作为先修：所需反对称算子、平方根、相容性等式全部收录。标准基中的 $J_0$ 约定为 $e_i\mapsto f_i$、$f_i\mapsto-e_i$；这来自 Lecture 7 的紧前例子。主目标使用 Lecture 8 Corollary 1 的完整证明，不扩大到原讲义接着仅以机制说明的全空间可缩性；也不纳入没有证明的 Floer/Arnold 结论。原 Lecture 7 第2页和 Lecture 8 第1页图像已读取。难度依据是构造在切丛上自然且光滑的结构及联结路径，而不是看到矩阵就称为谱方法。
\subsection*{可修改的简短标注（非作者正文）}
$c$：辛流形上相容几乎复结构的存在与变形。

$a$：辛配对与度量的同构；反对称算子；极分解；正定平方根；切丛逐点构造；度量空间的凸性。

\texttt{cross\_theory\_status = yes}。把几何相容性转成反对称算子的极分解，并利用度量的凸插值返回几乎复结构的路径。位置：Lecture 7 Proposition 1 与 Lecture 8 Corollary 1。

全部标注均为非穷尽、可修改初稿；未标注不表示未使用。
\subsection*{许可与修改记录}
材料来自 MIT OpenCourseWare，作者与课程见来源。按该站所示 CC BY-NC-SA 4.0 许可复用，整理部分沿用同一许可。\url{https://creativecommons.org/licenses/by-nc-sa/4.0/}。2026-09-28：助手转录题证、调整数学排版并加入来源、范围和初稿标注；不暗示作者或 MIT 认可本次整理。所留为本次题证转录摘录，不是原 PDF 下载件。
\end{document}
